\chapter{Giới thiệu chung}\label{chap:introduction}

Chương này xác định bài toán Hệ thống gợi ý hội thoại (CRS) và các thách thức còn tồn tại, tổng quan các hướng nghiên cứu liên quan, chỉ ra khoảng trống nghiên cứu, và trình bày mục tiêu cùng đóng góp kỹ thuật của khóa luận.


\section{Đặt vấn đề và Động lực nghiên cứu}

Khi khối lượng nội dung số bùng nổ và các nền tảng giải trí, thương mại điện tử ngày càng đa dạng, các hệ thống gợi ý truyền thống (Recommender Systems - RS) đang đối mặt với giới hạn cấu trúc khó vượt qua. Các mô hình lọc cộng tác và dựa trên nội dung vốn chỉ khai thác lịch sử tương tác một chiều (click, view, purchase), không thể nắm bắt sở thích tức thời hay các nhu cầu ngữ cảnh phức tạp phát sinh trong thực tế \cite{li2018towards, christakopoulou2016towardscrs}. Hệ thống gợi ý hội thoại (Conversational Recommender System - CRS) ra đời để lấp đầy khoảng trống này: thông qua giao tiếp bằng ngôn ngữ tự nhiên, CRS cho phép hệ thống chủ động làm rõ nhu cầu của người dùng và cung cấp các gợi ý có tính thuyết phục qua nhiều lượt tương tác.

Sự trưởng thành của các Mô hình ngôn ngữ lớn (Large Language Models - LLMs) mở ra cơ hội lớn cho CRS nhờ năng lực hiểu và sinh ngôn ngữ tự nhiên vượt trội. Tuy nhiên, ứng dụng trực tiếp LLM vào CRS bộc lộ hai rào cản cốt lõi. Về \textit{tính tin cậy}, LLM sinh văn bản theo phân phối xác suất chứ không truy xuất dữ liệu thực tế, khiến chúng dễ phát sinh hiện tượng "ảo giác" (hallucination), tự bịa ra các thực thể không tồn tại như tên phim, đạo diễn hay năm phát hành. Về \textit{hiệu quả vận hành}, việc tinh chỉnh LLM cho mỗi thay đổi dữ liệu tốn kém tài nguyên; ngược lại, nếu đưa toàn bộ tri thức vào context window thì chi phí API và độ trễ tăng vọt, ảnh hưởng trực tiếp đến khả năng phản hồi thời gian thực.

Để giải quyết bài toán ảo giác mà không cần huấn luyện lại mô hình, các kiến trúc kết hợp RAG (Retrieval-Augmented Generation) với Đồ thị tri thức (Knowledge Graph - KG) đã được đề xuất. Tiêu biểu trong số này là kiến trúc GRACE (Graph-Reasoning Augmented Conversational Engine) \cite{phan2025grace}, nghiên cứu tiền đề của chính tác giả đã được công bố tại SOICT 2025, vận hành bằng cách đối soát (grounding) kết quả sinh bởi LLM với dữ liệu thực trong KG thông qua cơ chế Truy xuất lai. Mặc dù giải quyết hiệu quả vấn đề ảo giác, GRACE và các hệ thống RAG tương tự vẫn hoạt động theo kiến trúc đường ống tĩnh (Static Pipeline), bộc lộ hạn chế rõ rệt khi đối mặt với các truy vấn yêu cầu suy luận đa bước hoặc chứa các ràng buộc mâu thuẫn, tình huống phổ biến trong hội thoại thực tế.

Điểm mấu chốt bị thiếu trong các hệ thống RAG tĩnh là \textit{cơ chế phản hồi nội bộ}: khi một bước xử lý thất bại (truy xuất rỗng, thực thể sai), toàn bộ pipeline bị tắc mà không có khả năng tự điều chỉnh. Đây là động lực để khóa luận đề xuất hệ thống \textbf{ARGOS} (Adaptive Reasoning and Graph Orchestration System), một khung làm việc đa tác tử có khả năng tự phát hiện điểm nghẽn, điều chỉnh chiến lược truy xuất và duy trì mạch hội thoại ngay cả trong các tình huống phức tạp.

\section{Tổng quan các nghiên cứu liên quan}

Nghiên cứu về Hệ thống Gợi ý Hội thoại đã trải qua ba thế hệ tiếp cận rõ rệt, mỗi thế hệ giải quyết hạn chế của thế hệ trước nhưng đồng thời mang theo những ràng buộc mới.

\subsection{Hệ thống CRS truyền thống tăng cường tri thức}
Các hệ thống CRS thế hệ đầu mô hình hóa luồng hội thoại bằng mạng nơ-ron hồi quy (RNN) để trích xuất sở thích và thực hiện gợi ý. Mô hình nền tảng ReDial \cite{li2018towards} kết hợp bộ tự mã hóa (autoencoder) với phân tích cảm xúc, song thường sinh ra các phản hồi "an toàn", chung chung và thiếu cá nhân hóa.

Nhận thấy hạn chế này, các nghiên cứu tiếp theo tích hợp Đồ thị Tri thức (KG) làm nguồn dữ liệu bổ trợ:
\begin{itemize}
    \item \textbf{KBRD \cite{chen-etal-2019-towards}:} Ánh xạ thực thể trong hội thoại vào đồ thị DBpedia và dùng GNN để lan truyền thông tin qua các quan hệ liên kết, cải thiện đáng kể độ chính xác so với ReDial.
    \item \textbf{KGSF \cite{zhou2020improvingcrs}:} Giải quyết sự mất kết nối giữa không gian từ vựng và không gian thực thể bằng cơ chế hợp nhất thông tin tương hỗ (Mutual Information Maximization), đồng bộ hóa biểu diễn giữa KG và mạng đồ thị từ vựng.
    \item \textbf{CR-Walker \cite{ma2021crwalker} và KERL \cite{Qiu_2025}:} Chuyển từ tính toán vector thuần túy sang mô hình hóa suy luận như chuỗi duyệt cây trên đồ thị. KERL còn sử dụng Học tăng cường (Reinforcement Learning) để tìm đường đi tối ưu trên KG.
\end{itemize}
Hạn chế cốt lõi của thế hệ này nằm ở module NLU cứng nhắc: phụ thuộc vào thuật toán nhận diện thực thể theo quy tắc khớp chuỗi (exact match), dễ bỏ sót thông tin khi người dùng diễn đạt ẩn ý hoặc dùng cách nói không chuẩn. Ngoài ra, mỗi tập dữ liệu mới đòi huấn luyện lại từ đầu với chi phí tài nguyên lớn.

\subsection{Mô hình Ngôn ngữ Lớn trong Hệ thống gợi ý}
LLM (GPT-4, LLaMA, Gemini) thay thế toàn bộ các module NLU và NLG truyền thống, cho phép hiểu ngữ cảnh phức tạp mà không cần quy tắc thủ công. Ba hướng tiếp cận chính được phân loại theo mức độ can thiệp vào trọng số mô hình:
\begin{itemize}
    \item \textbf{Zero/Few-Shot Prompting:} LLM-ReDial \cite{liang-etal-2024-llm} và LLM-ZCR \cite{he2023largelanguagemodels} giữ nguyên trọng số (frozen weights), chỉ thiết kế prompt. Triển khai nhanh nhưng hoàn toàn phụ thuộc vào trí nhớ tham số của LLM, dẫn đến ảo giác nghiêm trọng khi domain knowledge không được cập nhật.
    \item \textbf{Fine-tuning:} TallRec \cite{bao2023tallrec} đóng gói bài toán thành phân loại nhị phân; ReFICR \cite{yang-2024-reficr} dùng instruction tuning. Cả hai cải thiện độ chính xác nhưng kém linh hoạt khi có item mới, và chi phí tái huấn luyện liên tục là rào cản thực tiễn \cite{meng2025balancingfinetuning}.
    \item \textbf{Knowledge-Augmented Prompting:} UniCRS \cite{Wang_2022}, MESE \cite{yang-etal-2022-improving} và PECRS \cite{ravaut-etal-2024-parameter} tiêm tri thức đồ thị vào context của LLM. Giải pháp trung dung nhưng bị giới hạn bởi context window, chỉ có thể trích xuất đồ thị ở độ sâu 1-hop, đánh mất khả năng suy luận đa bước.
\end{itemize}

\subsection{Các kiến trúc lai và tăng cường truy xuất (RAG)}
RAG biến LLM từ "người học thuộc lòng" thành "người tra cứu tài liệu theo thời gian thực", vừa giảm ảo giác vừa tránh được chi phí fine-tuning:
\begin{itemize}
    \item \textbf{DCRS \cite{dao2024broadening}:} Dùng bộ truy xuất tương phản (Contrastive Retriever) để ánh xạ tín hiệu cộng tác (item-item) vào không gian vector của LLM, song vẫn gặp khó khăn với các truy vấn ngôn ngữ tự do.
    \item \textbf{COLA \cite{Lin_Wang_Li_2023}:} Cải thiện tích hợp cộng tác từ giai đoạn huấn luyện sớm, nhưng phạm vi chỉ xoay quanh ReDial.
    \item \textbf{G-CRS \cite{qiu2025graphragllm}:} Tiên phong dùng RAG trực tiếp trên KG: truy vấn ý định người dùng để quét KG, trích xuất mạng con (subgraph) liên quan rồi giao LLM đánh giá.
\end{itemize}

Điểm nghẽn chung của các kiến trúc RAG là khâu re-ranking: đưa hàng chục ứng viên kèm metadata vào Generative LLM làm tăng vọt độ trễ và chi phí API. Kiến trúc GRACE \cite{phan2025grace} đã giải quyết điểm này bằng cách kết hợp Hybrid Retrieval với cơ chế re-ranking dùng Generative LLM theo lô (batching), đạt hiệu năng vượt trội. Tuy nhiên, toàn bộ các hệ thống kể trên, kể cả GRACE, đều hoạt động theo đường ống tĩnh: dữ liệu chỉ chảy một chiều và không có cơ chế xử lý khi pipeline thất bại giữa chừng.

\section{Khoảng trống nghiên cứu và Sự cần thiết của hệ thống đa tác tử}

Bảng \ref{tab:related_work_summary} hệ thống hóa các hướng tiếp cận CRS theo bốn chiều so sánh. Nhìn vào bức tranh tổng thể, một hạn chế cấu trúc xuyên suốt tất cả các thế hệ: \textbf{tất cả đều vận hành theo cơ chế đường ống tĩnh (Static Pipeline)}.

\begin{table}[h!]
    \centering
    \caption{Tổng hợp các nghiên cứu liên quan trong lĩnh vực Hệ thống gợi ý hội thoại.}
    \label{tab:related_work_summary}
    \small
    \renewcommand{\arraystretch}{1.5}
    \begin{tabular}{@{}
        >{\raggedright\arraybackslash}p{2.6cm}
        >{\raggedright\arraybackslash}p{3.2cm}
        >{\raggedright\arraybackslash}p{4.6cm}
        >{\raggedright\arraybackslash}p{1.8cm}
        >{\raggedright\arraybackslash}p{1.6cm}
        @{}}
        \toprule
        \textbf{Phân loại} & \textbf{Mô hình tiêu biểu}                                                                             & \textbf{Phương pháp cốt lõi}                                      & \textbf{Dataset}          & \textbf{Cần huấn luyện} \\
        \midrule

        \multirow{3}{=}{\textbf{Truyền thống \& Tích hợp KG}}
                           & ReDial \cite{li2018towards}                                                                            & Bộ tự mã hóa RNN kết hợp phân tích cảm xúc.                       & ReDial                    & Có                      \\
        \cmidrule{2-5}
                           & KBRD \cite{chen-etal-2019-towards}, KGSF \cite{zhou2020improvingcrs}                                   & Tích hợp KG làm giàu biểu diễn người dùng.                        & ReDial, INSPIRED          & Có                      \\
        \cmidrule{2-5}
                           & CR-Walker \cite{ma2021crwalker}, KERL \cite{Qiu_2025}                                                  & Suy luận trên cấu trúc cây của KG và văn bản.                     & ReDial, INSPIRED          & Có                      \\
        \midrule

        \multirow{3}{=}{\textbf{LLM (Prompt/ Fine-tune)}}
                           & LLM-ZCR \cite{he2023largelanguagemodels}, LLM-ReDial \cite{liang-etal-2024-llm}                        & Prompt zero-shot/few-shot giữ đóng băng trọng số LLM.             & ReDial                    & Không                   \\
        \cmidrule{2-5}
                           & TallRec \cite{bao2023tallrec}, ReFICR \cite{yang-2024-reficr}                                          & Tinh chỉnh toàn bộ trọng số (Fine-tune).                          & Nhiều                     & Có                      \\
        \cmidrule{2-5}
                           & UniCRS \cite{Wang_2022}, MESE \cite{yang-etal-2022-improving}, PECRS \cite{ravaut-etal-2024-parameter} & Tăng cường kiến thức qua trích xuất KG tĩnh.                      & ReDial, INSPIRED          & Linh hoạt               \\
        \midrule

        \multirow{4}{=}{\textbf{Kiến trúc Lai \& RAG}}
                           & DCRS \cite{dao2024broadening}                                                                          & Bộ truy xuất tương phản (Contrastive Retriever).                  & ReDial, INSPIRED          & Có                      \\
        \cmidrule{2-5}
                           & COLA \cite{Lin_Wang_Li_2023}                                                                           & Tăng cường cộng tác huấn luyện từ đầu.                            & ReDial                    & Có                      \\
        \cmidrule{2-5}
                           & G-CRS \cite{qiu2025graphragllm}                                                                        & Graph RAG kết hợp với LLM nền tảng.                               & ReDial, INSPIRED          & Không                   \\
        \cmidrule{2-5}
                           & \textbf{GRACE (Baseline)}                                                                              & \textbf{Truy xuất lai (Sem+Content+Graph) kết hợp SLM Re-ranker.} & \textbf{ReDial, INSPIRED} & \textbf{Không}          \\
        \bottomrule
    \end{tabular}
\end{table}

Trong các pipeline tĩnh, dữ liệu chảy một chiều: phân tích ý định $\rightarrow$ truy xuất KG $\rightarrow$ xếp hạng $\rightarrow$ sinh phản hồi. Khi bất kỳ bước nào gặp sự cố (truy xuất không trả kết quả vì điều kiện quá khắt khe, hoặc LLM trích xuất nhầm thực thể), toàn bộ pipeline tắc mà không có cơ chế tự điều chỉnh. Kết quả là hệ thống hoặc trả về tập rỗng, hoặc sinh gợi ý chất lượng thấp, làm gián đoạn mạch hội thoại.

Phân tích kỹ hơn, hai khoảng trống nghiên cứu cụ thể được xác định:

\textbf{(1) Thiếu cơ chế tự kiểm định (Self-Correction):} Không có hệ thống CRS hiện tại nào có khả năng phát hiện và sửa lỗi nội bộ trong quá trình xử lý. Ví dụ cụ thể: nếu LLM trích xuất sai tên thực thể và gợi ý không tồn tại trong KG, hệ thống không thể nhận biết và loại bỏ ứng viên đó trước khi trả về người dùng.

\textbf{(2) Thiếu khả năng thỏa hiệp ràng buộc (Constraint Relaxation):} Khi người dùng đặt ra các điều kiện mâu thuẫn hoặc quá hạn hẹp, các hệ thống tĩnh không thể tự quyết định nên nới lỏng ràng buộc nào để đưa ra gợi ý xấp xỉ; thay vào đó chúng trả về kết quả rỗng, một trải nghiệm người dùng rất kém.

Trong khi đó, nghiên cứu về kiến trúc đa tác tử dựa trên LLM \cite{xi2023rise} đã chứng minh rằng việc phân tách hệ thống thành các tác tử chuyên biệt với vòng lặp phản hồi (feedback loop) có thể xử lý hiệu quả các bài toán đòi hỏi suy luận đa bước và tự sửa lỗi. Tuy nhiên, hướng tiếp cận này \textit{chưa được áp dụng vào bài toán CRS kết hợp Đồ thị tri thức}, đây chính là khoảng trống mà khóa luận đề xuất lấp đầy thông qua hệ thống \textbf{ARGOS}.

\section{Mục tiêu nghiên cứu}

Khóa luận đặt ra bốn mục tiêu nghiên cứu:
\begin{itemize}
    \item \textbf{Mục tiêu 1:} Phân tích và định lượng các điểm nghẽn của kiến trúc RAG tĩnh trong hệ thống CRS, cụ thể qua mô hình GRACE, để tạo cơ sở thực nghiệm cho việc đề xuất cơ chế mới.
    \item \textbf{Mục tiêu 2:} Thiết kế kiến trúc đa tác tử ARGOS với cơ chế phân phối nhiệm vụ và vòng lặp phản hồi nội bộ, giải quyết hạn chế của pipeline tĩnh mà không yêu cầu huấn luyện lại mô hình.
    \item \textbf{Mục tiêu 3:} Phát triển và tích hợp cơ chế Critic-Relaxation (bao gồm Critic Agent kiểm định tính xác thực và Relaxation Agent nới lỏng ràng buộc) vào vòng lặp suy luận của ARGOS.
    \item \textbf{Mục tiêu 4:} Đánh giá toàn diện ARGOS trên hai bộ dữ liệu chuẩn ReDial và INSPIRED, so sánh với các mô hình tiên tiến nhất hiện nay, và phân tích định lượng đóng góp của từng thành phần.
\end{itemize}

\section{Đóng góp của khóa luận}

Khóa luận mang lại ba đóng góp kỹ thuật mới cho lĩnh vực CRS:
\begin{enumerate}
    \item \textbf{Kiến trúc đa tác tử ARGOS:} Lần đầu tiên áp dụng mô hình Multi-Agent với vòng lặp suy luận-hành động vào hệ thống CRS kết hợp Đồ thị tri thức. ARGOS phân tách pipeline thành năm tác tử chuyên biệt (Profiler Agent, Orchestrator, Graph Agent, Critic Agent, Relaxation Agent) và sử dụng Dynamic Weighting để điều phối chiến lược truy xuất theo hình thái ngôn ngữ của truy vấn, cơ chế chưa xuất hiện trong các hệ thống CRS prior work.
    \item \textbf{Cơ chế Critic-Relaxation:} Phương pháp tự kiểm định và thỏa hiệp ràng buộc tự động trong CRS. Critic Agent đối chiếu trực tiếp ứng viên với KG Neo4j để loại bỏ hallucination; Relaxation Agent áp dụng thuật toán phân cấp ưu tiên (Sacrifice Hierarchy) để tái cấu trúc truy vấn khi gặp tập kết quả rỗng. Đây là đóng góp kỹ thuật độc lập, có thể tách rời và ứng dụng cho các hệ thống RAG khác ngoài CRS.
    \item \textbf{Kết quả thực nghiệm và mã nguồn mở:} ARGOS cải thiện đáng kể so với kiến trúc tiền đề GRACE và đạt kết quả cạnh tranh với các hệ thống tiên tiến trên ReDial và INSPIRED về chỉ số Recall@K. Toàn bộ mã nguồn và dữ liệu thực nghiệm được công bố công khai tại \url{https://github.com/DatPhan06/GRACE}, tạo nền tảng cho các nghiên cứu tiếp theo trong lĩnh vực.
\end{enumerate}

\section*{Kết luận chương}

Chương 1 đã xác lập nền tảng định hướng cho toàn bộ khóa luận. Xuất phát từ hạn chế cốt lõi của các hệ thống CRS hiện tại (thiếu tính linh hoạt của thế hệ truyền thống và hiện tượng ảo giác của các LLM thuần túy), phân tích cho thấy các kiến trúc RAG, dù đã cải thiện đáng kể, vẫn bị ràng buộc bởi cơ chế đường ống tĩnh. Khoảng trống về khả năng tự kiểm định và thỏa hiệp ràng buộc là cơ sở kỹ thuật rõ ràng để đề xuất hệ thống ARGOS. Các kiến thức nền tảng cần thiết để hiểu kiến trúc ARGOS sẽ được trình bày chi tiết trong Chương 2.
